\documentclass[10pt,a4paper]{amsart}
\usepackage{amsmath,amssymb,amsthm,graphicx}
\usepackage{verbatim}
\newenvironment{reference}{\comment}{\endcomment}
%\newenvironment{reference}{}{}
\newenvironment{slogan}{\comment}{\endcomment}
\newenvironment{history}{\comment}{\endcomment}

% For commutative diagrams we use Xy-pic
\usepackage[all]{xy}

% We use 2cell for 2-commutative diagrams.
\xyoption{2cell}
% Theorem environments.
%
\theoremstyle{plain}
\newtheorem{theorem}[subsection]{Theorem}
\newtheorem{proposition}[subsection]{Proposition}
\newtheorem{lemma}[subsection]{Lemma}

\theoremstyle{definition}
\newtheorem{definition}[subsection]{Definition}
\newtheorem{example}[subsection]{Example}
\newtheorem{exercise}[subsection]{Exercise}
\newtheorem{situation}[subsection]{Situation}

\theoremstyle{remark}
\newtheorem{remark}[subsection]{Remark}
\newtheorem{remarks}[subsection]{Remarks}

\numberwithin{equation}{subsection}

% Macros
%
\def\lim{\mathop{\mathrm{lim}}\nolimits}
\def\colim{\mathop{\mathrm{colim}}\nolimits}
\def\Spec{\mathop{\mathrm{Spec}}}
\def\Hom{\mathop{\mathrm{Hom}}\nolimits}
\def\Ext{\mathop{\mathrm{Ext}}\nolimits}
\def\SheafHom{\mathop{\mathcal{H}\!\mathit{om}}\nolimits}
\def\SheafExt{\mathop{\mathcal{E}\!\mathit{xt}}\nolimits}
\def\Sch{\mathit{Sch}}
\def\Mor{\mathop{\mathrm{Mor}}\nolimits}
\def\Ob{\mathop{\mathrm{Ob}}\nolimits}
\def\Sh{\mathop{\mathit{Sh}}\nolimits}
\def\NL{\mathop{N\!L}\nolimits}
\def\CH{\mathop{\mathrm{CH}}\nolimits}
\def\proetale{{pro\text{-}\acute{e}tale}}
\def\etale{{\acute{e}tale}}
\def\QCoh{\mathit{QCoh}}
\def\Ker{\mathop{\mathrm{Ker}}}
\def\Im{\mathop{\mathrm{Im}}}
\def\Coker{\mathop{\mathrm{Coker}}}
\def\Coim{\mathop{\mathrm{Coim}}}

% Boxtimes
%
\DeclareMathSymbol{\boxtimes}{\mathbin}{AMSa}{"02}

%
% Macros for moduli stacks/spaces
%
\def\QCohstack{\mathcal{QC}\!\mathit{oh}}
\def\Cohstack{\mathcal{C}\!\mathit{oh}}
\def\Spacesstack{\mathcal{S}\!\mathit{paces}}
\def\Quotfunctor{\mathrm{Quot}}
\def\Hilbfunctor{\mathrm{Hilb}}
\def\Curvesstack{\mathcal{C}\!\mathit{urves}}
\def\Polarizedstack{\mathcal{P}\!\mathit{olarized}}
\def\Complexesstack{\mathcal{C}\!\mathit{omplexes}}
% \Pic is the operator that assigns to X its picard group, usage \Pic(X)
% \Picardstack_{X/B} denotes the Picard stack of X over B
% \Picardfunctor_{X/B} denotes the Picard functor of X over B
\def\Pic{\mathop{\mathrm{Pic}}\nolimits}
\def\Picardstack{\mathcal{P}\!\mathit{ic}}
\def\Picardfunctor{\mathrm{Pic}}
\def\Deformationcategory{\mathcal{D}\!\mathit{ef}}
\usepackage{cleveref}
\crefname{originallemma}{original lemma}{original lemmas}
\Crefname{originallemma}{Original lemma}{Original lemmas}
\makeatletter
\newlabel{more-morphisms-lemma-case-of-tor-independence}{{0CTA}{source}{}{sourcepointer-1}{}}
\newlabel{more-morphisms-lemma-case-of-tor-independence@cref}{{[originallemma][0][]0CTA}{[1][0][]source}{}{}{}}
\makeatother
\crefname{originallemma}{original lemma}{original lemmas}
\Crefname{originallemma}{Original lemma}{Original lemmas}
\makeatletter
\newlabel{cohomology-lemma-derived-pullback-composition}{{0D5S}{source}{}{sourcepointer-2}{}}
\newlabel{cohomology-lemma-derived-pullback-composition@cref}{{[originallemma][0][]0D5S}{[1][0][]source}{}{}{}}
\makeatother
\crefname{originallemma}{original lemma}{original lemmas}
\Crefname{originallemma}{Original lemma}{Original lemmas}
\makeatletter
\newlabel{perfect-lemma-base-change-relatively-perfect}{{0DI5}{source}{}{sourcepointer-3}{}}
\newlabel{perfect-lemma-base-change-relatively-perfect@cref}{{[originallemma][0][]0DI5}{[1][0][]source}{}{}{}}
\makeatother
\crefname{originallemma}{original lemma}{original lemmas}
\Crefname{originallemma}{Original lemma}{Original lemmas}
\makeatletter
\newlabel{lemma-relative-dualizing-complex-algebra}{{0E2U}{source}{}{sourcepointer-4}{}}
\newlabel{lemma-relative-dualizing-complex-algebra@cref}{{[originallemma][0][]0E2U}{[1][0][]source}{}{}{}}
\makeatother
\crefname{originallemma}{original lemma}{original lemmas}
\Crefname{originallemma}{Original lemma}{Original lemmas}
\makeatletter
\newlabel{dualizing-lemma-base-change-relative-dualizing}{{0E2F}{source}{}{sourcepointer-5}{}}
\newlabel{dualizing-lemma-base-change-relative-dualizing@cref}{{[originallemma][0][]0E2F}{[1][0][]source}{}{}{}}
\makeatother
\title{A relative dualizing complex pulls back to a relative dualizing complex}
\author{The Stacks Project Authors}
\date{stacks-project:duality.tex:local-snapshot-20261012; GNU FDL 1.2 or later}
\begin{document}
\maketitle
\noindent\textit{Editorial note.} All delivered mathematical text is the author's own, copied from the pinned Stacks Project chapter file duality.tex; no statement or proof step was rewritten, reordered or removed.\par\medskip
\noindent\textit{Editorial note.} Preamble provenance: the pinned chapter file opens with an input of the shared preamble and carries no preamble text of its own, so the project's single Stacks-wide preamble staged at sources/stacks-duality/preamble.tex is excerpted verbatim under the recipe fixed by LANE\_BRIEF 22.1, as three spans, lines 8-13 (verbatim package and the reference/slogan/history environments), lines 14-18 (Xy-pic setup) and lines 154-217 (the native theorem declarations and mathematical macros). Nothing was added, removed or redefined; the wrapper class and the amsmath/amssymb/amsthm/graphicx base packages are supplied by the assembler.\par\medskip
\noindent\textit{Editorial note.} The counted statement and proof contain no cross reference, no citation and no included file, so the delivered body is self contained.\par\medskip
\noindent\textit{Editorial note.} Canonical location of the counted proposition: Stacks tag 0E2Y (sources/stacks-duality/pinned-original-source.tex line 7471).\par\medskip
\noindent\textit{Allowed standard prerequisites (editorial).}\begin{itemize}
\item Relative dualizing complexes and their description through the diagonal. Same chapter, the definition of a relative dualizing complex preceding this lemma, invoked in the statement, and Lemma labelled lemma-relative-dualizing-complex-algebra, cited at proof lines 7518-7520. Fixes the data and turns the affine local check into algebra.
\item Tor independence of the diagonal situations and the resulting base change of the diagonal. More on Morphisms, Lemma labelled more-morphisms-lemma-case-of-tor-independence, cited at proof lines 7496-7499. Gives L(g' x g')\textasciicircum{}* Delta\_*O\_X = Delta'\_*O\_\{X'\}.
\item Transitivity of derived pullback and relative perfectness under base change. Cohomology, Lemma labelled cohomology-lemma-derived-pullback-composition, cited at proof lines 7499-7502, and Derived Categories of Schemes, Lemma labelled perfect-lemma-base-change-relatively-perfect, cited at proof lines 7508-7510. Identifies the pulled back comparison map and shows K' is S'-perfect.
\item Base change of the algebraic relative dualizing complex. Dualizing Complexes, Lemma labelled dualizing-lemma-base-change-relative-dualizing, cited at proof lines 7520-7524. Supplies the algebraic statement the proof reduces to.
\end{itemize}\medskip
\ifdefined\hypertarget\hypertarget{sourcepointer-1}{}\fi
\noindent\textit{Source reference.} Original lemma 0CTA: sources/stacks-chow/tags.txt, source lines 16060--16060. Cross-chapter standard prerequisite located in the project tag list staged in the local snapshot; invoked at proof line 7508 of the counted proof. It is retained only as a navigation pointer; every load-bearing step is inside the delivered span.\par\medskip
\ifdefined\hypertarget\hypertarget{sourcepointer-2}{}\fi
\noindent\textit{Source reference.} Original lemma 0D5S: sources/stacks-chow/tags.txt, source lines 16497--16497. Cross-chapter standard prerequisite located in the project tag list staged in the local snapshot; invoked at proof line 7510 of the counted proof. It is retained only as a navigation pointer; every load-bearing step is inside the delivered span.\par\medskip
\ifdefined\hypertarget\hypertarget{sourcepointer-3}{}\fi
\noindent\textit{Source reference.} Original lemma 0DI5: sources/stacks-chow/tags.txt, source lines 16930--16930. Cross-chapter standard prerequisite located in the project tag list staged in the local snapshot; invoked at proof line 7517 of the counted proof. It is retained only as a navigation pointer; every load-bearing step is inside the delivered span.\par\medskip
\ifdefined\hypertarget\hypertarget{sourcepointer-4}{}\fi
\noindent\textit{Source reference.} Original lemma 0E2U: sources/stacks-duality/pinned-original-source.tex, source lines 7271--7271. Same-chapter pointer into the author environment lemma at source line 7271; invoked at proof line 7523 of the counted proof. It is retained only as a navigation pointer; every load-bearing step is inside the delivered span.\par\medskip
\ifdefined\hypertarget\hypertarget{sourcepointer-5}{}\fi
\noindent\textit{Source reference.} Original lemma 0E2F: sources/stacks-chow/tags.txt, source lines 17605--17605. Cross-chapter standard prerequisite located in the project tag list staged in the local snapshot; invoked at proof line 7526 of the counted proof. It is retained only as a navigation pointer; every load-bearing step is inside the delivered span.\par\medskip
\begin{lemma}
\label{lemma-base-change-relative-dualizing}
Consider a cartesian square
$$
\xymatrix{
X' \ar[d]_{f'} \ar[r]_{g'} & X \ar[d]^f \\
S' \ar[r]^g & S
}
$$
of schemes. Assume $X \to S$ is flat and locally of finite presentation.
Let $(K, \xi)$ be a relative dualizing complex for $f$.
Set $K' = L(g')^*K$. Let $\xi'$ be the derived base change of $\xi$
(see proof). Then $(K', \xi')$ is a relative dualizing complex for $f'$.
\end{lemma}

\begin{proof}
Consider the cartesian square
$$
\xymatrix{
X' \ar[d]_{\Delta_{X'/S'}} \ar[r] & X \ar[d]^{\Delta_{X/S}} \\
X' \times_{S'} X' \ar[r]^{g' \times g'} & X \times_S X
}
$$
Choose $W \subset X \times_S X$ open such that $\Delta_{X/S}$
factors through a closed immersion $\Delta : X \to W$.
Choose $W' \subset X' \times_{S'} X'$ open such that $\Delta_{X'/S'}$
factors through a closed immersion $\Delta' : X \to W'$
and such that $(g' \times g')(W') \subset W$. Let us still denote
$g' \times g' : W' \to W$ the induced morphism. We have
$$
L(g' \times g')^*\Delta_*\mathcal{O}_X =
\Delta'_*\mathcal{O}_{X'}
\quad\text{and}\quad
L(g' \times g')^*L\text{pr}_1^*K|_W =
L\text{pr}_1^*K'|_{W'}
$$
The first equality holds because $X$ and $X' \times_{S'} X'$
are tor independent over $X \times_S X$ (see for example
More on Morphisms, Lemma \ref{more-morphisms-lemma-case-of-tor-independence}).
The second holds by transitivity of derived pullback
(Cohomology, Lemma \ref{cohomology-lemma-derived-pullback-composition}).
Thus $\xi' = L(g' \times g')^*\xi$ can be viewed as a map
$$
\xi' : \Delta'_*\mathcal{O}_{X'} \longrightarrow L\text{pr}_1^*K'|_{W'}
$$
Having said this the proof of the lemma is straightforward.
First, $K'$ is $S'$-perfect by Derived Categories of Schemes, Lemma
\ref{perfect-lemma-base-change-relatively-perfect}.
To check that $\xi'$ induces an isomorphism
of $\Delta'_*\mathcal{O}_{X'}$ to
$R\SheafHom_{\mathcal{O}_{W'}}(
\Delta'_*\mathcal{O}_{X'}, L\text{pr}_1^*K'|_{W'})$
we may work affine locally. By
Lemma \ref{lemma-relative-dualizing-complex-algebra}
we reduce to the corresponding statement in algebra
which is proven in Dualizing Complexes, Lemma
\ref{dualizing-lemma-base-change-relative-dualizing}.
\end{proof}

\end{document}
